% Generated draft: requires mathematical and source review.
\documentclass[11pt]{article}
\usepackage[margin=25mm]{geometry}
\usepackage{fontspec}
\setmainfont{Latin Modern Roman}
\usepackage{amsmath,amssymb,mathtools,unicode-math}
\setmathfont{Latin Modern Math}
\usepackage{enumitem,needspace,xurl,hyperref}
\setlength{\parindent}{0pt}
\setlength{\parskip}{5pt}
\setlength{\emergencystretch}{3em}
\newcommand{\N}{\mathbb{N}}
\newcommand{\Z}{\mathbb{Z}}
\newcommand{\Q}{\mathbb{Q}}
\newcommand{\R}{\mathbb{R}}
\newcommand{\C}{\mathbb{C}}
\newcommand{\F}{\mathbb{F}}
\newcommand{\A}{\mathbb{A}}
\newcommand{\PP}{\mathbb{P}}
\newcommand{\E}{\mathbb{E}}
\newcommand{\eps}{\varepsilon}
\newcommand{\dd}{\,\mathrm{d}}
\begin{document}
\section*{UnsolvedMath ID: 30002451}
\textbf{Source problem code:} catalogue entry
\subsection{Definitions and mathematical statement}
Let \(G\) be a simple complex Lie group and let \(\mathfrak g=\operatorname{Lie}(G)\). Choose Cartan and simple-root data for \(\mathfrak g\), and denote the set of simple roots by \(\Delta\). A Belavin--Drinfeld triple is a triple \(T=(\Gamma_1,\Gamma_2,\gamma)\), where \(\Gamma_1,\Gamma_2\subseteq\Delta\), the bijection \(\gamma:\Gamma_1\to\Gamma_2\) preserves the root inner product, and for every \(\alpha\in\Gamma_1\) there exists \(k\geq1\) such that \(\gamma^j(\alpha)\in\Gamma_1\) for \(0\leq j<k\), while \(\gamma^k(\alpha)\notin\Gamma_1\). Thus iteration eventually leaves the domain of \(\gamma\). Let \(R_T\) denote the associated Belavin--Drinfeld class of classical \(r\)-matrices, and write \(\{\cdot,\cdot\}_r\) for the associated Poisson--Lie bracket on \(G\). The supplied record does not specify the normalization or representative convention for \(R_T\).

A cluster structure on \(G\) is understood as a realization of a cluster algebra by rational functions on \(G\), with each extended cluster forming a transcendence basis of the rational function field \(\C(G)\) and seeds related by cluster mutation. An extended cluster includes mutable variables and frozen variables. Compatibility with \(\{\cdot,\cdot\}_r\) means that every extended cluster \((x_1,\ldots,x_m)\) is log-canonical for that bracket: there are constants \(\omega_{ij}^{r,\Sigma}\in\C\), depending on \(r\) and the seed \(\Sigma\), such that
\[
\{x_i,x_j\}_r=\omega_{ij}^{r,\Sigma}x_ix_j
\qquad(1\leq i,j\leq m).
\]
The upper cluster algebra \(\mathcal U(\mathcal C_T)\) is the intersection, inside \(\C(G)\), of the Laurent rings in the mutable variables of all seeds, with coefficients in the selected frozen-variable coefficient ring. Whether frozen variables are inverted must be specified; the supplied record does not specify this convention. A cluster structure is regular if all its cluster variables are regular functions on \(G\).

The existence question in the supplied clean statement is: for every simple complex Lie group \(G\) and every Belavin--Drinfeld triple \(T\) for its Lie algebra, does there exist a single cluster structure \(\mathcal C_T\) such that
\[
\begin{gathered}
\text{for every }r\in R_T,\quad
\mathcal C_T\text{ is compatible with }\{\cdot,\cdot\}_r,\\
\mathcal U(\mathcal C_T)\cong\mathcal O(G),
\end{gathered}
\]
where the isomorphism is the natural one induced by realizing the cluster variables as functions on \(G\)? In particular, the same cluster structure must work for every \(r\in R_T\).

There are unresolved qualifications to this formulation. Interpreting \(\mathcal O(G)\) as an algebraic coordinate ring requires an algebraic form of \(G\); the supplied phrase ``simple complex Lie group'' does not specify the global form or distinguish algebraic regular functions from holomorphic functions. The displayed question therefore retains the intended notation \(\mathcal O(G)\) without selecting an unstated global form or coefficient convention.

The supplied original statement instead asks whether there is a classification of regular cluster structures on \(G\) completely parallel to the Belavin--Drinfeld classification. It does not specify the equivalence relation on cluster structures, the precise correspondence, or the meaning of completeness. Existence of the structures displayed above does not by itself specify such a classification. The original regularity qualification and the distinction between these two targets require mathematical review.
\subsection{Short English statement}
For every simple complex Lie group \(G\) and Belavin--Drinfeld triple \(T\), the clean statement asks for one cluster structure compatible with every Poisson bracket arising from \(R_T\), whose upper cluster algebra naturally recovers \(\mathcal O(G)\). The original statement asks more broadly for a parallel classification of regular cluster structures. The relationship between these targets, the global form of \(G\), and the frozen-variable convention remain unspecified. The supplied status is partially solved; this is a dated curation label, not an independently verified conclusion.
\subsection{Sources}
\begin{itemize}
\item Supplied local UnsolvedMath registry record, permanent ID \(30002451\), title ``Cluster Structures for Belavin--Drinfeld Data,'' source problem code \texttt{catalogue entry}, in the collection ``Oberwolfach Reports Open Problems.'' The record gives proposed year 2013 and names no individual proposer.
\item Supplied report citation: \emph{Cluster Algebras and Related Topics (2014)}, Oberwolfach Reports, pp.~3379--3432, DOI \texttt{10.4171/owr/2013/58}. \url{https://doi.org/10.4171/owr/2013/58}. This is also the supplied literature source labeled ``original OWR report.'' The supplied background describes these reports as documenting mathematical workshops and their open-problem discussions.
\item Supplied literature source labeled ``Gekhtman Shapiro and Vainshtein'': \url{https://arxiv.org/abs/1902.02902}.
\item Supplied literature source labeled ``Zhang SO3 cluster structures'': \url{https://doi.org/10.1080/00927872.2025.2450036}.
\item Literature assessment dated \texttt{2026-08-21}: the supplied curation labels the problem partially solved and reports verification for standard data and broad classes of type A data, with isolated further examples, while describing the all-groups, all-triples classification as remaining open. This is supplied dated curation derived from the OWR record, not an independently verified literature assessment. Its time-sensitive status and cited claims require independent verification.
\item The supplied statement status is \texttt{corrected\_verified}, with a supplied assertion that the clean statement was checked against the official Oberwolfach report and preserves its hypotheses, notation, and open target. That assertion is recorded as source metadata; no independent source verification is claimed here.
\end{itemize}
\end{document}
